\documentclass[9pt,letterpaper]{extarticle}
\usepackage[T1]{fontenc}
\usepackage[utf8]{inputenc}
% Font: \timestrue = Times, \timesfalse = Latin Modern (font of the reference CV)
\newif\iftimes
\timestrue
\iftimes
  \usepackage{mathptmx}
\else
  \usepackage{lmodern}
\fi
\usepackage[letterpaper,top=0.55in,bottom=0.6in,left=0.62in,right=0.62in]{geometry}
\usepackage{xcolor,enumitem,titlesec,tabularx,array,microtype}
\usepackage[hidelinks]{hyperref}
\definecolor{link}{HTML}{2F6FC4}
\hypersetup{colorlinks=true,urlcolor=link,linkcolor=link}
\pagestyle{empty}
\setlength{\parindent}{0pt}
\setlength{\tabcolsep}{0pt}
\titleformat{\section}{\large\bfseries\color{link}}{}{0pt}{}[{\vspace{-2pt}\color{black}\titlerule[0.4pt]}]
\titlespacing*{\section}{0pt}{7pt}{3pt}
\setlist[itemize]{leftmargin=1.2em,itemsep=0pt,topsep=3pt,parsep=0pt,label=\textbullet}
\newcommand{\me}{\underline{\textbf{Kyulee Jeon}}}
\newcommand{\mes}{\underline{\textbf{Kyulee Jeon}$^*$}}
\newcommand{\role}[1]{(\textit{#1})}
\newcommand{\head}[2]{\textbf{#1}\hfill\textit{#2}\par}
\newcommand{\subhead}[2]{#1\hfill #2\par}
\newcolumntype{L}[1]{>{\raggedright\arraybackslash}p{#1}}
\newcolumntype{Y}{>{\raggedright\arraybackslash}X}
\newenvironment{rows}[1]{\tabularx{\textwidth}{@{}L{#1}@{\hspace{0.12in}}Y@{}}}{\endtabularx}
\newcommand{\gap}{\\\noalign{\vskip 2.5pt}}
\newcommand{\lgap}{\\\noalign{\vskip 4pt}}
\newcommand{\sep}{\quad$|$\quad}

\begin{document}
\begin{center}
{\huge Kyulee Jeon}\\[3pt]
Ph.D. Candidate in Biomedical Systems Informatics, Yonsei University College of Medicine\\[1pt]
\href{mailto:jklee320@yuhs.ac}{jklee320@yuhs.ac}\sep
\href{https://kyulee-jeon.github.io}{kyulee-jeon.github.io}\sep
\href{https://scholar.google.com/citations?user=WhhJ_usAAAAJ}{Google Scholar}\sep
\href{https://github.com/kyulee-jeon}{GitHub}\sep
\href{https://www.linkedin.com/in/kyulee-jeon-870415202}{LinkedIn}
\end{center}
\vspace{-4pt}

\section{Research Profile}
Ph.D. candidate in Biomedical Systems Informatics at Yonsei University College of Medicine, advised by Prof. Seng Chan You.
I make \textbf{multimodal health data interoperable} so that clinical AI can be validated and adopted in routine care.
My research links multimodal medical images with electronic health records through a standardized common data model,
measures the real-world data heterogeneity that limits data reuse and AI deployment, and generates evidence for clinical AI
through multi-institutional, distributed validation. Expected graduation: August 2027.\par\vspace{2pt}
\textbf{Research interests:} multimodal health data interoperability; medical imaging data standardization; common data models;
DICOM metadata; evidence generation for clinical artificial intelligence; distributed AI validation; observational research;
terminology standardization with large language models.

\section{Research Directions}
\begin{rows}{1.25in}
\textbf{Multimodal data interoperability} & Standardized common data models that link multimodal medical images with electronic health records for the secondary use of clinical data.\gap
\textbf{Imaging data standardization} & Real-world DICOM metadata heterogeneity and the essential metadata that reproducible research and AI deployment need.\gap
\textbf{Evidence for clinical AI} & Multi-institutional, distributed validation pipelines that give the evidence to adopt AI models in routine clinical care.\gap
\textbf{Terminology standardization} & Ontology-driven methods, with large language models, that map local clinical terms to international standards.
\end{rows}

\section{Featured Papers}
\textbf{Development of Medical Imaging Data Standardization for Imaging-Based Observational Research: OMOP Common Data Model Extension} \role{co-first author}\par
Woo Yeon Park$^*$, \mes, Teri Sippel Schmidt, Haridimos Kondylakis, Tarik Alkasab, Blake E. Dewey, Seng Chan You, Paul Nagy\par
\textit{Journal of Imaging Informatics in Medicine}, 2024 \ \href{https://doi.org/10.1007/s10278-024-00982-6}{paper}
\begin{itemize}
\item Proposes the Medical Imaging CDM, the imaging extension for OMOP CDM: two new tables and two vocabularies that link imaging events and image-derived features with standardized clinical data.
\item The basis of my later work on multimodal data ETL and AI validation, and of the work of the OHDSI Medical Imaging Working Group.
\end{itemize}
\vspace{3pt}
\textbf{Development of Clinically Validated Artificial Intelligence Model for Detecting ST-segment Elevation Myocardial Infarction} \role{co-first author}\par
Sang-Hyup Lee$^*$, \mes, Yong-Joon Lee, Seng Chan You, Seung-Jun Lee, Sung-Jin Hong, Chul-Min Ahn, Jung-Sun Kim, Byeong-Keuk Kim, Young-Guk Ko, Donghoon Choi, Myeong-Ki Hong\par
\textit{Annals of Emergency Medicine}, 2024 \ \href{https://doi.org/10.1016/j.annemergmed.2024.06.004}{paper}
\begin{itemize}
\item \textbf{Clinically validated} AI model: developed with ECGs of a percutaneous coronary intervention cohort to reflect the clinical context of STEMI patients, and validated with ECGs of emergency room patients with chest pain.
\item Supervised learning with labels from two cardiologists who reviewed coronary angiography.
\item Confirmed the benefit of the model in clinical settings: it identified additional true negatives for FIRST CP activation.
\end{itemize}
\vspace{3pt}
\textbf{ARKE: An Ontology-Driven Framework for Automated Mapping of Local Radiology Procedure Terms to the LOINC-RadLex Playbook Using Large Language Model} \role{co-first author}\par
Sumin Lee$^*$, \mes, Joon Seok Lim, Yiju Park, Min Seong Kim, Ju Hyun Jin, Chang Hoon Han, Soon Ho Yoon, Seng Chan You\par
\textit{Journal of Biomedical Informatics}, 2026 \ \href{https://doi.org/10.1016/j.jbi.2026.105071}{paper}
\begin{itemize}
\item \textbf{Why it matters}: multi-institutional validation of imaging AI needs the exam name that links imaging exams to patient data, but each hospital writes it differently. Manual mapping to the LOINC/RSNA Radiology Playbook is slow and inconsistent between readers.
\item \textbf{Mapping as knowledge encoding}: ARKE decomposes each local name into Playbook attributes, retrieves the top 10 candidates from the ontology by rule-based scoring, and confines the LLM to attribute parsing and final selection. No unconstrained generation, so hallucination is minimized and no codes are invented.
\item \textbf{Result}: across two tertiary hospitals in Korea (3,326 procedure names), recall@1 of 0.70 (internal) and 0.81 (external), more than 20 points above direct code-name mapping (0.48 and 0.50; p < 0.001).
\end{itemize}

\section{Selected Research Projects}
\head{Multimodal Data Interoperability: Imaging Extension for OMOP CDM (MI-CDM)}{2022--present}
\subhead{Yonsei University, with Johns Hopkins University and the OHDSI Medical Imaging Working Group}{}
\begin{itemize}
\item \textbf{MI-CDM}: co-authored the imaging extension for OMOP CDM. Two extension tables link imaging events and image-derived features with clinical data.
\item \textbf{MIMIC-ECG to MI-CDM}: ETL of MIMIC-IV ECG data into the extension tables. \href{https://github.com/kyulee-jeon/MIMIC-ECG}{github.com/kyulee-jeon/MIMIC-ECG}
\item \textbf{MIMIC-CXR to MI-CDM}: ETL of MIMIC chest X-ray data into the extension tables. \href{https://github.com/kyulee-jeon/MIMIC-CXR}{github.com/kyulee-jeon/MIMIC-CXR}
\item \textbf{End-to-end AI validation pipeline} on this data model: case study with chest X-rays of tuberculosis patients (OHDSI Global Symposium 2026).
\item \textbf{Multi-institutional validation network} (\textit{project manager}; funded by a Seoul National University donation fund): built the same data structure with Seoul National University Hospital and SMG-SNU Boramae Medical Center, and tested the AI validation pipeline across the institutions.
\end{itemize}
\pagebreak
\head{Real-World DICOM Metadata Heterogeneity}{2024--present}
\subhead{Yonsei University}{}
\begin{itemize}
\item \textbf{Nationwide network study} (\textit{first author}, IEEE Journal of Biomedical and Health Informatics, major revision): characterizes DICOM metadata heterogeneity in radiology across institutions.
\item \textbf{DICOM compliance assessment}: compliance with mandatory DICOM attributes. Excellence Prize, KOSMI 2024.
\end{itemize}
\vspace{2pt}
\head{Clinical AI Development and Terminology Standardization}{2022--present}
\subhead{Yonsei University}{}
\begin{itemize}
\item \textbf{ARKE} (\textit{co-first author}, J Biomed Inform 2026): ontology-driven LLM framework that maps local radiology procedure names to the LOINC/RSNA Radiology Playbook.
\item \textbf{STEMI detection} (\textit{co-first author}, Ann Emerg Med 2024): deep ensemble ECG model with clinical validation. Excellence Prize, KOSMI 2022.
\end{itemize}

\section{Education}
\head{Yonsei University College of Medicine}{Sep 2022--Aug 2027 (expected)}
\subhead{Integrated M.S./Ph.D., Biomedical Systems Informatics; advisor: Prof. Seng Chan You}{Seoul, Korea}
\vspace{2pt}
\head{Ewha Womans University}{Mar 2017--Feb 2021}
\subhead{B.S., Public Health; GPA 4.0/4.5}{Seoul, Korea}

\section{Publications}
$^*$ denotes equal contribution.\par\vspace{2pt}
\begin{rows}{1.0in}
\textbf{IEEE JBHI\newline Major revision} & \textbf{Nationwide Characterization of Real-World DICOM Metadata Heterogeneity in Radiology: A Network Cross-Sectional Study} \role{first author}\lgap
\textbf{JMIR Med Inform\newline In press} & \textbf{Integrating Pathology Images and Reports Into the OMOP Common Data Model: A Practical Framework Using DICOM Metadata} \role{co-author}\lgap
\textbf{J Imaging Inform Med 2026} & \textbf{Advancing the FAIRness of Multimodal Imaging Research Through the OMOP MI-CDM Framework: A Case Replication Study in Alzheimer's Disease} \role{co-author}\newline Gabriel L. O. Salvador, Jen Wooyeon Park, Teri Sippel Schmidt, Yiju Park, \me, Seng Chan You, Paul Nagy, Blake Dewey, Alzheimer's Disease Neuroimaging Initiative\newline \href{https://doi.org/10.1007/s10278-026-02146-0}{paper}\lgap
\textbf{J Biomed Inform 2026} & \textbf{ARKE: An Ontology-Driven Framework for Automated Mapping of Local Radiology Procedure Terms to the LOINC-RadLex Playbook Using Large Language Model} \role{co-first author}\newline Sumin Lee$^*$, \mes, Joon Seok Lim, Yiju Park, Min Seong Kim, Ju Hyun Jin, Chang Hoon Han, Soon Ho Yoon, Seng Chan You\newline \href{https://doi.org/10.1016/j.jbi.2026.105071}{paper}\lgap
\textbf{J Evid-Based Pract 2026} & \textbf{Evaluating Regional Diversity in Scientific Communication: A Comparative Analysis of COVID-19 Preprints and Peer-Reviewed Publications} \role{co-first author}\newline Dong Hyun Kim$^*$, \mes, Seng Chan You\newline \href{https://doi.org/10.63528/jebp.2026.00007}{paper}\lgap
\textbf{JAMIA 2025} & \textbf{Breaking Data Silos: Incorporating the DICOM Imaging Standard Into the OMOP CDM to Enable Multimodal Research} \role{co-author}\newline Woo Yeon Park, Teri Sippel Schmidt, Gabriel Salvador, Kevin O'Donnell, Brad Genereaux, \me, Seng Chan You, Blake E. Dewey, Paul Nagy, Alzheimer's Disease Neuroimaging Initiative\newline \href{https://doi.org/10.1093/jamia/ocaf091}{paper}\lgap
\textbf{Invest Radiol 2025} & \textbf{Advancing Medical Imaging Research Through Standardization: The Path to Rapid Development, Rigorous Validation, and Robust Reproducibility} \role{first author}\newline \me, Woo Yeon Park, Charles E. Kahn Jr., Paul Nagy, Seng Chan You, Soon Ho Yoon\newline \href{https://doi.org/10.1097/RLI.0000000000001106}{paper}\lgap
\textbf{Ann Emerg Med 2024} & \textbf{Development of Clinically Validated Artificial Intelligence Model for Detecting ST-segment Elevation Myocardial Infarction} \role{co-first author}\newline Sang-Hyup Lee$^*$, \mes, Yong-Joon Lee, Seng Chan You, Seung-Jun Lee, Sung-Jin Hong, Chul-Min Ahn, Jung-Sun Kim, Byeong-Keuk Kim, Young-Guk Ko, Donghoon Choi, Myeong-Ki Hong\newline \href{https://doi.org/10.1016/j.annemergmed.2024.06.004}{paper}\lgap
\textbf{J Imaging Inform Med 2024} & \textbf{Development of Medical Imaging Data Standardization for Imaging-Based Observational Research: OMOP Common Data Model Extension} \role{co-first author}\newline Woo Yeon Park$^*$, \mes, Teri Sippel Schmidt, Haridimos Kondylakis, Tarik Alkasab, Blake E. Dewey, Seng Chan You, Paul Nagy\newline \href{https://doi.org/10.1007/s10278-024-00982-6}{paper}
\end{rows}

\newpage
\section{Conference Presentations}
\begin{rows}{1.0in}
Nov 2026 & \textbf{OHDSI APAC Symposium}, Research in Action talk (main conference). Medical Imaging CDM in Action: Toward an Imaging Evidence Network.\gap
Nov 2026 & \textbf{OHDSI APAC Symposium}, poster. Building a Longitudinal Multimodal Dataset Using MI-CDM: Standardizing MIMIC-IV ECG and Chest Radiography Data.\gap
Aug 2026 & \textbf{KCR}, poster. Implementation of Distributed AI Model Validation Framework Based on OMOP Medical Image CDM: A Case Study Using Chest X-rays of Tuberculosis Patients.
\end{rows}

\vspace{2.5pt}
\begin{rows}{1.0in}
Oct 2025 & \textbf{OHDSI Global Symposium}, poster. Implementation of Distributed AI Model Validation Framework Based on OMOP Medical Image CDM.\gap
Jun 2025 & \textbf{KOSMI}, poster. Implementation of Distributed AI Model Validation Framework Based on OMOP Medical Image CDM.\gap
Jun 2025 & \textbf{KOSMI}, poster. ARKE: An Automated System for Standardizing Radiology Procedure Terminology Using LLMs and RAG.
\end{rows}

\vspace{2.5pt}
\begin{rows}{1.0in}
Oct 2024 & \textbf{OHDSI APAC Symposium}, lightning talk and poster. Feasibility of Integrating DICOM Headers into the OMOP Medical Imaging CDM: A Pilot Study Using Chest CT Data.
\end{rows}

\vspace{2.5pt}
\begin{rows}{1.0in}
Jul 2024 & \textbf{K-SIIM}, poster. Development of Medical Imaging Data Standardization for Imaging Based Observational Research: OMOP CDM Extension.\gap
Jun 2024 & \textbf{KOSMI}, poster. Assessing Compliance with Mandatory Metadata Attributes in DICOM, the International Standard for Medical Imaging. \textbf{Excellence Prize}.\gap
Nov 2023 & \textbf{KOSMI}, poster. Extension of the OMOP Common Data Model for Medical Imaging Data Standardization.
\end{rows}

\vspace{2.5pt}
\begin{rows}{1.0in}
Oct 2023 & \textbf{OHDSI Global Symposium}, poster. Development of Medical Imaging Data Standardization for Imaging Based Observational Research: OMOP CDM Extension.\gap
Jun 2023 & \textbf{KOSMI}, poster. Development of Medical Imaging Data Standardization for Imaging Based Observational Research: OMOP CDM Extension.\gap
Nov 2022 & \textbf{KOSMI}, oral. Development of STEMI Detection Algorithm with 12-lead ECG Signals and 1D CNN. \textbf{Excellence Prize}.
\end{rows}

\section{Awards and Service}
\begin{rows}{1.0in}
Oct 2025 & \textbf{1st place}. OHDSI Early-Stage Researchers Working Group.\gap
Jun 2024 & \textbf{Excellence Prize for Poster Presentation}. Korean Society of Medical Informatics (KOSMI) Spring Conference.\gap
Nov 2022 & \textbf{Excellence Prize for Oral Presentation}. KOSMI Autumn Conference.\gap
Oct 2020 & \textbf{Grand Prize for Poster Presentation}. Korean Society for Health Education and Promotion Autumn Conference.\gap
Nov 2018 & \textbf{Grand Prize}. Video on responsible drinking habits, Korea Health Promotion Institute.\gap
Service & \textbf{Reviewer}. npj Digital Medicine (4), Scientific Reports (2), Scientific Data (2), BMC Medical Imaging (1).\gap
Nov 2026 & \textbf{Lead organizer}. 2026 OHDSI APAC Datathon: prepared standardized multimodal data, designed and ran the program, and coordinated the organizing team.
\end{rows}

\section{Teaching}
\begin{rows}{1.0in}
2026 & Tutorial instructor, Medical Imaging Working Group Tutorial (ETL), OHDSI Global Symposium.\gap
2026 & Organizer and instructor, advanced track, Yonsei University College of Medicine Data Bootcamp.\gap
2026 & Data organizer and instructor, Yonsei University Health System Challenge.\gap
Nov 2024 & Instructor, ECG data pre-challenge training, Yonsei University Health System Big Data Challenge.\gap
May 2022 & Mini-tutorial on concept sets and cohort definition, OHDSI APAC Community Call.
\end{rows}

\section{Technical Skills}
\begin{rows}{1.25in}
\textbf{Data standards} & OMOP CDM and its imaging extension (MI-CDM), DICOM, LOINC/RSNA Radiology Playbook, RadLex, SNOMED CT, OHDSI tools (ATLAS, cohort definition).\gap
\textbf{Data engineering} & ETL pipelines for imaging, waveform, and clinical data; SQL; DICOM metadata extraction and quality assessment.\gap
\textbf{Clinical AI} & Deep learning for ECG and medical imaging, distributed model validation, LLM and retrieval-augmented generation pipelines.\gap
\textbf{Programming} & Python, SQL, R, Git.\gap
\textbf{Certifications} & Health Information Manager and Health Educator (2021), Ministry of Health and Welfare, Korea.
\end{rows}
\end{document}
